% Generated draft: requires mathematical and source review.
\documentclass[11pt]{article}
\usepackage[margin=25mm]{geometry}
\usepackage{fontspec}
\setmainfont{Latin Modern Roman}
\usepackage{amsmath,amssymb,mathtools,unicode-math}
\setmathfont{Latin Modern Math}
\usepackage{enumitem,needspace,xurl,hyperref}
\setlength{\parindent}{0pt}
\setlength{\parskip}{5pt}
\setlength{\emergencystretch}{3em}
\newcommand{\N}{\mathbb{N}}
\newcommand{\Z}{\mathbb{Z}}
\newcommand{\Q}{\mathbb{Q}}
\newcommand{\R}{\mathbb{R}}
\newcommand{\C}{\mathbb{C}}
\newcommand{\F}{\mathbb{F}}
\newcommand{\A}{\mathbb{A}}
\newcommand{\PP}{\mathbb{P}}
\newcommand{\E}{\mathbb{E}}
\newcommand{\eps}{\varepsilon}
\newcommand{\dd}{\,\mathrm{d}}
\begin{document}
\section*{UnsolvedMath ID: 30004742}
\textbf{Source problem code:} catalogue entry
\subsection{Definitions and mathematical statement}
\textbf{Source conjecture.} Euclidean billiards have no absolutely periodic orbits.

The supplied record does not define ``absolutely periodic'' or specify the dimension, boundary regularity, convexity, or admissible class of billiard tables. Consequently, a fully specified universal conjecture cannot be recovered without further source clarification. The following gives an explicit conditional interpretation consistent with the record's distinction between finite-order degeneracies and infinite-order examples.

In a Euclidean billiard table \(Q\subset\R^d\), a trajectory travels along straight segments and reflects at a regular boundary collision according to
\[
 v_{\mathrm{out}}=v_{\mathrm{in}}-2\langle v_{\mathrm{in}},n(x)\rangle n(x),
\]
where \(n(x)\) is a unit normal at the collision point \(x\). Let \(T\) denote the collision map on collision states consisting of a boundary point and an outgoing unit velocity. Its iterates are considered only where the relevant collisions and return map are defined. For the interpretation below, restrict to non-tangent collisions at smooth boundary points for which the required iterate is \(C^\infty\) on a neighborhood of the state.

A periodic state satisfies \(T^m(p)=p\) for some integer \(m\geq1\); here \(m\) need not be the least period. If ``absolutely periodic'' means that an iterate agrees with the identity to infinite order at a periodic state, its precise local condition is as follows. In local coordinates \(\varphi\) near \(p\), put \(z_0=\varphi(p)\) and
\[
 F(z)=\varphi\bigl(T^m(\varphi^{-1}(z))\bigr)-z.
\]
The condition is
\[
 D^\alpha F(z_0)=0
 \qquad\text{for every multi-index }\alpha\in\N_0^r,
\]
where \(r\) is the dimension of the collision phase space and \(\N_0=\{0,1,2,\ldots\}\). This includes the periodicity condition when \(|\alpha|=0\).

Under this interpretation, the conjecture asserts that for every table in the intended admissible class, every integer \(m\geq1\), and every regular periodic state \(p\) with a smooth local return map, there exists a multi-index \(\alpha\) with \(|\alpha|\geq1\) such that
\[
 D^\alpha F(z_0)\ne0.
\]
The intended admissible class and the identification of absolute periodicity with this infinite-order condition require confirmation. Vanishing through any prescribed finite order does not by itself establish vanishing at every order.
\subsection{Short English statement}
The source conjectures that Euclidean billiards have no absolutely periodic trajectories. Interpreting this as infinite-order agreement of a local return map with the identity gives the conditional statement above. The intended billiard class and terminology remain unspecified; the supplied partial-resolution assessment is dated curation requiring independent review.
\subsection{Sources}
\begin{itemize}
\item Supplied UnsolvedMath local registry record, permanent ID \(30004742\), exact source problem code \texttt{catalogue entry}, in the Oberwolfach Reports Open Problems collection. The record gives proposed year 2021 and names no proposer.
\item Supplied report citation: \emph{Dynamische Systeme (2022)}, Oberwolfach Reports, DOI 10.4171/owr/2021/33, pp.\ 1735--1803. \url{https://doi.org/10.4171/owr/2021/33}. The supplied background describes Oberwolfach Reports as documenting mathematical workshops and their open-problem discussions.
\item Supplied literature reference, labeled ``arbitrarily high finite-order orbits'': \url{https://arxiv.org/abs/2209.11721}.
\item Supplied literature reference, labeled ``analytic billiard result'': \url{https://arxiv.org/abs/math/0110134}.
\item Literature assessment dated exactly \texttt{2026-08-22}: supplied dated curation, not independently verified. It labels the problem partially solved and says that the no-absolutely-periodic-orbits conjecture remains unresolved; analytic convex cases and periods through four are ruled out, while arbitrarily high finite-order absolutely periodic degeneracies are constructed, rather than infinite-order examples. The supplied background identifies this as literature triage derived from the supplied OWR curation and requests independent verification of time-sensitive status and cited claims.
\item The registry labels the statement \texttt{corrected\_verified} and reports manual semantic, typographic, and LaTeX review. These are supplied editorial assertions, not independent verification of the formulation or mathematics here.
\end{itemize}
\end{document}
